\documentclass{article}
\usepackage{natbib}
\usepackage{graphicx}
\def\method{FooNet\xspace}
\newcommand{\data}{BarSet}
\newcommand{\todo}[1]{{\color{red}#1}}
% \def\unused{should not appear}

\title{A Fixture Paper}
\author{Ada Lovelace \and Alan Turing}

\begin{document}
\maketitle

\begin{abstract}
We study \method on the \data benchmark.
\end{abstract}

\section{Introduction}
Deep networks dominate vision benchmarks~\cite{smith2020,jones2019}.
% A comment line inside a paragraph must not split it.
However, they are data hungry \cite{lee2021}. We address this.


\section{Method}
\label{sec:method}
As shown in Fig.~\ref{fig:arch}, \method follows prior work. % trailing comment
The model reaches $12\%$ higher accuracy than the baseline~\cite{smith2020}.
\begin{figure}[t]
  \includegraphics[width=0.8\linewidth]{arch.pdf}
  \caption{A figure caption citing \cite{jones2019}, which must be ignored.}
  \label{fig:arch}
\end{figure}
\citet{lee2021} report that pretraining helps in low-resource settings.
Training is stable for learning rates up to $10^{-3}$ (see~\citealp[Sec.~3]{smith2020}).

\iffalse
This paragraph is disabled \cite{jones2019}.
\fi

\section{Related Work}
Several methods exist, e.g.\ adapters and prompts \citep{jones2019, lee2021}. Dr.\ Smith et al.\ proposed the first version.
Unrelated sentence without citations.
Early work studied this extensively. \cite{missingkey}
Results improved further.\footnote{See also \cite{jones2019}.}

\bibliographystyle{plain}
\bibliography{refs}
\end{document}
